\section{Méthodologie et Plateformes de travail}

\begin{frame}{Topologies Matérielles}
    \begin{itemize}
        \item Tests effectués sur deux plateformes aux caractéristiques distinctes :
    \end{itemize}
    \begin{table}
        \footnotesize
        \centering
        \begin{tabular}{|l|c|c|}
            \hline
            \textbf{Caractéristique} & \textbf{Plateforme AMD} & \textbf{Plateforme Intel} \\ \hline
            Processeur & Ryzen 7 7840HS & Core i5-13500H \\ \hline
            Cœurs / Threads & 8 cœurs / 16 threads & 12 cœurs / 16 threads \\ \hline
            Fréquence Max & 3801 MHz & 4700 MHz \\ \hline
            Cache L2 & 8 MiB (Dédié) & 9 MiB (Partagé) \\ \hline
            Cache L3 & 16 MiB (Partagé) & 18 MiB (Partagé) \\ \hline
            Topologie & 1 nœud NUMA & 1 nœud NUMA \\ \hline
        \end{tabular}
    \end{table}
   \begin{block}{Idée}
        Réaliser une expérience comparative pour observer comment l'architecture influence le programme selon le paradigme de langage ou le compilateur utilisé.
    \end{block}
\end{frame}

\begin{frame}{Paradigmes de Langage}
    \begin{itemize}
        \item \textbf{Fortran 90} : Benchmark officiel NPB MG de la NASA.
        \item \textbf{C++} : Version disponible sur GitHub.
    \end{itemize}
    \begin{block}{Pourquoi le Fortran ?}
        \begin{itemize}
            \item \textbf{Par curiosité} : Langage historique du calcul scientifique, nous voulions voir s'il restait plus performant que le C++.
            \item \textbf{Évaluation de MAQAO} : Tester l'outil sur un autre langage classique (Fortran) face à un langage objet (C++).
            \item \textbf{Gestion des tableaux} : Vérifier si l'optimisation native du Fortran facilite le travail du compilateur.
        \end{itemize}
    \end{block} 
\end{frame}

\begin{frame}{Stratégie de Compilation et Automatisation}
    \begin{itemize}
        \item \textbf{Classe de test} : Sélection de la \textbf{Classe C} ($512^3$ points) pour saturer les caches.
        \item \textbf{Toolchains} : Solutions constructeurs (AOCC, Intel oneAPI) et base commune (GCC, Clang++).
        \item \textbf{Optimisations} : \texttt{-O3}, \texttt{-march=native}, \texttt{-ffast-math}, \texttt{-funroll-loops}.
    \end{itemize}
    \begin{block}{Automatisation}
        Mise en place de scripts pour automatiser la génération des nombreux rapports MAQAO et gérer les combinaisons de langages.
    \end{block}
\end{frame}

\begin{frame}{Protocole de Mesure avec MAQAO}
     Le script génère plusieurs types de rapports pour chaque binaire :
    \begin{description}
        \item[Analyse de base (-R1)] Localisation des fonctions les plus coûteuses.
        \item[Stabilité (-S1)] \textbf{10 répétitions} pour garantir la fiabilité de nos mesures face aux variations du système.
        \item[Scalabilité (-WS)] Étude de la montée en charge (Strong Scaling) de 1 à 8 cœurs physiques.
        \item[Duels différentiels] Comparaison automatique des binaires (ex: C++ vs Fortran) pour isoler les écarts de performance.
    \end{description}
\end{frame}

\begin{frame}{Fiabilisation de l'Environnement}
    \begin{itemize}
        \item \textbf{Affinité} : Utilisation de \texttt{export OMP\_PLACES=cores} pour empêcher l'OS de déplacer les threads.
        \item \textbf{Fréquence} : Fixation en mode \textbf{performance} avant les tests (\texttt{sudo cpupower frequency-set -g performance}).
    \end{itemize}
    \begin{block}{Indicateur de préparation}
        Ces réglages ont permis de stabiliser l'affinité à \textbf{77,9\%} (contre 23,4\% par défaut), garantissant la fiabilité des analyses.
    \end{block}
\end{frame}